%!TEX root = ../main.tex
% =============================================================================
%  一、问题重述与分析
%  评审在问：你知道题目在问什么吗？
%  本节含：① 记号约定与决策矛盾（式 (1)）  ② 四问的信息结构与新增数学对象（表 1）
%          ③ 决策时序图（图 1）、依赖图（图 2）  ④ 两条基线的口径分离（表 2、表 3）
%  写法约定：本节起全文用「集合／参数／决策变量／对偶变量」四类记号；
%            完整符号表见 §2.3，公式中每个符号都在其中登记。
% =============================================================================

\section{问题重述与分析}\label{chap:problem}

\subsection{记号与决策矛盾}

全天离散为 \(|\mathcal T|=144\) 个时段，\(t\in\mathcal T\)，步长 \(\Delta t=1/6\)~h；结果窗口为
\(\mathcal D\)，\(|\mathcal D|=334\) 天（2025-02-01 至 12-31）。第 \(t\) 槽的电价、负载功率与光伏功率记为
\(p_t,L_t,PV_t\)。储能充、放电量 \(c_t,d_t\) 与储电量 \(S_t\) 满足
\(0\le c_t,d_t\le E^{\max}=\overline P\Delta t\)、\(S_t\in[\underline S,\overline S]\)，
购电量非负 \(E_t\ge0\)（假设 A1：不外送）。

\textbf{决策矛盾。} 微网的收益只能来自电价的时间价差：把电量从低价时段搬到高价时段。储能搬运
\(1\)~kWh 完成一次循环须克服往返损耗，故套利有利可图的\textbf{盈亏平衡价比}为

\begin{equation}
  \frac{p_{\max}}{p_{\min}}\;>\;\frac{1}{\eta_c\eta_d}=1.2346 ,
  \label{eq:ch1-arb}
\end{equation}

而附件 1 的实测价比为 \(1.3952/0.3713=3.76\)，远高于该阈值——\textbf{价差不是瓶颈}。真正的约束是
容量摆幅、充放功率与两类不确定性（负载与光伏）。四问逐层引入这些约束：

\begin{itemize}
  \item \textbf{问题 1}（确定性计划）：\((p_t,L_t,PV_t)\) 已知且逐日重复，要求 \(S_0=S_{144}\)，
        在 \(0{:}00\) 一次制定全天计划；
  \item \textbf{问题 2}（日前计划）：\(p_t\) 仍固定，\((L_t,PV_t)\) 逐日随机；计划层在 \(0{:}00\) 决策，
        缺口在结算时刻按 \(5p_t\) 补足（\textbf{紧急购电}），构成两阶段决策；
  \item \textbf{问题 3}（滚动调整）：\(m\in\{0,6,12,18\}\) 时可获取未来 24~h 整点光伏预报，
        可在该时刻修正计划；计划高于调整的部分按 \(0.5p_t\) 计价、低于的部分按 \(1.5p_t\) 计价；
        追问是否需要引入其他时刻的预报；
  \item \textbf{问题 4}（波动电价）：\(p_t\) 亦随机，问题由“确定价格 \(+\) 随机出力”升级为
        \textbf{价格与出力的联合不确定性}，且须给出风险度量。
\end{itemize}

\subsection{一个模型族、两个开关}

四问共用同一套变量 \((E_t,c_t,d_t,S_t)\) 与同一组约束：电量平衡式~\eqref{eq:q1-balance}、
储电量递推式~\eqref{eq:q1-soc}、容量与功率界、以及假设 A1。它们仅由两个开关区分：
\textbf{电价向量}（附件 1 日内曲线／附件 4 实际序列）与\textbf{信息层数}
（1 个决策时刻／1 个决策时刻 \(+\) 情景／4 个决策时刻）。表~\ref{tab:fourq} 给出四问在该框架下的位置。

\begin{table}[htbp]
  \centering
  \caption{四问在“一个模型族、两个开关”下的位置（范围：问题 1 为附件 1 典型日，问题 2--4 为
           2025-02-01 至 12-31 共 334 天；单位与量纲见 \S\ref{sec:notation}；样本：4 个子问题）}
  \label{tab:fourq}
  \small
  \setlength{\tabcolsep}{4pt}
  \begin{tabularx}{\linewidth}{C{1.0cm} L{2.0cm} L{2.5cm} L{2.2cm} Y}
    \toprule
    子问题 & 电价向量 & 信息层数 & 目标结构 & 相对上一问\textbf{新增}的数学对象 \\
    \midrule
    1 & 附件 1 日内曲线 & 1 个决策时刻 & 确定性 LP & 对偶水价值 \(\nu_t\) 与阈值律（命题~\ref{prop:threshold}） \\
    2 & 附件 1 日内曲线 & 1 个决策时刻 \(+\) 情景 & 两阶段随机规划 & \textbf{情景维度}与追索量 \(g_{\omega,t}\)；报童临界分位数 \\
    3 & 附件 1 日内曲线 & 4 个决策时刻 & 多阶段（滚动实现） & \textbf{非预期性要求}式~\eqref{eq:q3-nonanticip}；信息价值 \\
    4 & 附件 4 实际序列 & 同上两种 & 同上 \(+\) 风险度量 & 电价随机过程、联合情景、\(\mathrm{CVaR}\) \\
    \bottomrule
  \end{tabularx}
\end{table}

末两列须强调：问题 2 相对问题 1 增加\textbf{情景维度}（计划量共享、追索量随情景变化），问题 3 相对
问题 2 增加\textbf{非预期性要求}与调整定价，二者均为结构性新增而非参数替换。决定建模方式的即三件事：
\textbf{决策时刻个数}、\textbf{每时刻的信息集}、\textbf{偏差的惩罚结构}（图~\ref{fig:timeline}）。
\textbf{一条全局声明}：问题 3（及问题 4-3）给出的多阶段模型\textbf{不展开场景树}，实际求解的是它的
滚动时域（MPC）可行策略，故该问的费用是\textbf{策略评价}而非多阶段最优值；理由、
边界与对信息价值结论的影响集中在 \S\ref{sec:q3-mpc}，并在摘要与\chapref{chap:eval}结论处重复标注。

\begin{figure}[htbp]
  \centering
  %!TEX root = ../main.tex
% =============================================================================
%  图 1  决策时序图（正文 §1.2）
%  规格（四项必须标注，缺一不可）：
%    ① "计划依据"与"结算依据"是两套数据，用不同线型区分
%    ② 18:00 的预报覆盖到次日 18:00，用一条跨越 24:00 的箭头表示
%    ③ 紧急购电发生在实时（结算时），不在决策时刻
%    ④ 四问在图上的差别：Q1 仅 0:00 一个决策点且无惩罚；Q2 有 0:00 + 实时惩罚；
%       Q3 有 4 个决策点；Q4 旁注"电价亦为随机"
% =============================================================================
\begin{tikzpicture}[
  x=0.46cm, y=1cm, font=\small, >=Stealth,
  ev/.style={font=\scriptsize, inner sep=1.5pt},
]

% ---------- 泳道底色 ----------
\fill[gray!8, rounded corners=2pt] (-3.15,-2.15) rectangle (26.2,-0.30);   % 信息
\fill[blue!5, rounded corners=2pt] (-3.15,-3.70) rectangle (26.2,-2.30);   % 决策
\fill[red!4,  rounded corners=2pt] (-3.15,-5.25) rectangle (26.2,-3.85);   % 结算

% ---------- 泳道标签 ----------
\node[rotate=90, font=\footnotesize\heiti, gray!60!black] at (-2.55,-1.22) {信息};
\node[rotate=90, font=\footnotesize\heiti, gray!60!black] at (-2.55,-3.00) {决策};
\node[rotate=90, font=\footnotesize\heiti, gray!60!black] at (-2.55,-4.55) {结算};

% ---------- 时间刻度（横轴 0:00 -> 24:00） ----------
\foreach \h in {0,6,12,18,24}{
  \draw[gray!25] (\h,-5.25) -- (\h,-0.30);
}
\foreach \h/\lab in {0/0:00, 6/6:00, 12/12:00, 18/18:00, 24/24:00}{
  \draw[gray!60] (\h,0.00) -- (\h,-0.16);
  \node[above, font=\scriptsize, gray!70!black] at (\h,0.00) {\lab};
}
\draw[gray!60, ->] (-0.9,0.00) -- (26.2,0.00) node[right, font=\scriptsize, gray!70!black] {时刻};
\draw[gray!45, dashed] (24,0.00) -- (24,-5.25);
\node[font=\scriptsize, gray!70!black, align=center] at (25.45,-5.55) {次日};

% ---------- 泳道 1：信息（附件 3 的四个预报发布时刻） ----------
% 预报覆盖区间的起点用实心圆标注
\foreach \h/\y/\len/\txt in {%
  0/-0.62/24/{\,覆盖至 24:00},
  6/-0.95/18/{\,覆盖至 24:00},
  12/-1.28/12/{\,覆盖至 24:00},
  18/-1.61/8/{~覆盖至次日 18:00（跨日）}}{
  \fill[blue!55!black] (\h,\y) circle (1.5pt);
  \draw[gray!60, line width=0.5pt, -{Stealth[length=4pt]}] (\h,\y) -- (\h+\len,\y);
  \node[ev, anchor=west, gray!25!black] at (\h+0.35,\y) {\h:00 发布\txt};
}

% ---------- 泳道 2：决策 ----------
% Q1：唯一决策时刻 0:00
\fill[blue!65!black] (0,-2.62) circle (2.4pt);
\node[font=\scriptsize, anchor=west, blue!50!black] at (0.35,-2.62)
  {计划购电：Q1\,/\,Q2\,/\,Q3\,/\,Q4（唯一决策时刻，随后不再修正）};
% Q3/Q4：6:00、12:00、18:00 调整
\foreach \h in {6,12,18}{
  \draw[blue!65!black, fill=white, line width=0.7pt] (\h,-3.28) circle (2.4pt);
}
\draw[gray!55, dotted, line width=0.6pt] (0,-2.62) -- (0,-3.28);
\node[font=\scriptsize, anchor=west, blue!50!black] at (6.35,-3.28)
  {调整购电：Q3\,/\,Q4 在 6:00、12:00、18:00 各修正一次};

% ---------- 泳道 3：结算 ----------
\draw[gray!55, densely dashed, line width=0.7pt, -{Stealth[length=5pt]}]
  (0,-4.20) -- (24,-4.20);
\node[font=\scriptsize, anchor=west, gray!25!black] at (0.35,-3.98)
  {结算依据：附件 2 实际负载与实际光伏功率（\textbf{与计划依据是两套数据}）};
% 紧急购电：发生在实时（结算时），不在决策时刻
\foreach \x in {9.4, 15.1, 21.3}{
  \draw[red!75!black, line width=0.9pt] (\x-0.28,-4.52) -- (\x+0.28,-4.52) -- (\x,-4.78) -- cycle;
}
\node[font=\scriptsize, anchor=west, red!65!black] at (-3.05,-4.98)
  {缺口 $\rightarrow$ 紧急购电（按 $5p_t$ 结算），发生在\textbf{实时}，不在任何决策时刻};

% ---------- 图例 ----------
\node[ev, anchor=west] at (-3.05,-5.62) {%
  \tikz\draw[gray!60,line width=0.5pt,-{Stealth[length=4pt]}] (0,0)--(0.55,0);
  ~预报（计划依据）\quad
  \tikz\draw[gray!55,densely dashed,line width=0.7pt,-{Stealth[length=4pt]}] (0,0)--(0.55,0);
  ~实际值（结算依据）\quad
  \tikz\fill[blue!65!black] (0,0) circle (2.2pt);~决策时刻\quad
  \tikz\draw[red!75!black,line width=0.9pt] (0,-0.05)--(0.28,-0.05)--(0.14,-0.26)--cycle;
  ~紧急购电（非决策时刻）\quad
  注：Q4 中电价亦为随机（旁注）};
\end{tikzpicture}

  \caption{微网决策时序与信息结构（上行泳道为信息到达时刻，中行为决策时刻，下行为结算依据；
            \(\Delta t=1/6\)~h，全天 144 个 10~min 时段；样本窗口 2025-02-01 至 12-31 共 334 天）}
  \label{fig:timeline}
\end{figure}

两侧惩罚方向相反，故不能用同一套边际规则处理。问题 3 中追加的边际代价是削减的边际收益的
\(1.5/0.5=3\) 倍，频繁改计划本身即付费；问题 2 与 4-2 则相反——计划量按正常电价全额结算，缺口按
\(5p_t\) 补足，故最优计划\textbf{主动多买}。消融实验（\S\ref{sec:ablation}）给出该结论的量级：
取消安全边际后 334 天总费用由 \(14\,264\,792.13\) 元升至 \(17\,398\,559.68\) 元（比无储能基线高
\(6.04\%\)），紧急购电量由 \(172\,015.82\)~kWh 升至 \(1\,156\,618.02\)~kWh，
边际价值 \(3\,133\,767.55\) 元。

\subsection{依赖关系与误差传播}

图~\ref{fig:dependency} 表明四问共享同一物理骨架（并网点电量平衡、储电量递推、充放电上下界），
逐层增加信息结构模块，使每层新增贡献可被\textbf{单独消融}（\chapref{chap:verify}）。

\begin{figure}[htbp]
  \centering
  %!TEX root = ../main.tex
% =============================================================================
%  图 2  四问的模型依赖关系（正文 §1.3）
%  规格：附件 -> Q1 -> Q2 -> Q3 -> Q4，每个箭头必须标注"提供什么"；
%        本图用于证明四问是一条信息结构递进链，而不是四篇互不相干的报告。
% =============================================================================
\begin{tikzpicture}[
  font=\small, >=Stealth,
  att/.style ={draw=gray!50, fill=gray!8, rounded corners=2pt, align=center,
               font=\scriptsize, text width=2.0cm, minimum height=0.72cm, inner sep=3pt},
  pre/.style ={draw=orange!55!black, fill=orange!8, rounded corners=2pt, align=center,
               font=\scriptsize, text width=6.6cm, minimum height=0.85cm, inner sep=4pt},
  qbox/.style={draw=blue!50, fill=blue!7, rounded corners=2pt, align=center,
               font=\small, text width=4.6cm, minimum height=0.95cm, inner sep=4pt},
  obox/.style={draw=green!45!black, fill=green!7, rounded corners=2pt, align=center,
               font=\scriptsize, text width=2.2cm, minimum height=0.8cm, inner sep=3pt},
  ar/.style  ={->, draw=gray!65, line width=0.6pt},
  lab/.style ={font=\scriptsize, anchor=west, align=left, text width=5.6cm,
               inner sep=1pt, gray!20!black},
]

% ---------------- 附件 ----------------
\node[att] (a1) at (-4.35,0) {附件 1\\电价/负载/光伏};
\node[att] (a2) at (-1.45,0) {附件 2\\负载/光伏实际值};
\node[att] (a3) at ( 1.45,0) {附件 3\\四个时刻光伏预报};
\node[att] (a4) at ( 4.35,0) {附件 4\\波动电价};

% ---------------- 数据预处理 ----------------
\node[pre] (pre) at (0,-1.35)
  {数据预处理：时标对齐（标签 $=$ 时段末端）\ \textbullet\ 附件 1 身份核查\ \textbullet\ 特征与误差分析（第三章）};

\foreach \n in {a1,a2,a3,a4}{\draw[ar] (\n.south) -- (pre.north);}

% ---------------- 四问 ----------------
\node[qbox] (q1) at (0,-2.95) {\textbf{问题 1}\ 确定性线性规划（LP）};
\node[qbox] (q2) at (0,-4.60) {\textbf{问题 2}\ 两阶段随机规划};
\node[qbox] (q3) at (0,-6.25) {\textbf{问题 3}\ 多阶段滚动 $+$ 非预期性};
\node[qbox] (q4) at (0,-7.90) {\textbf{问题 4}\ 波动电价 $+$ 风险度量};

\draw[ar] (pre.south) -- (q1.north);

% 箭头上的"提供什么"标注
\draw[ar] (q1.south) -- (q2.north);
\node[lab] at (0.35,-3.78)
  {\textbf{提供}\ 功率平衡与储电量递推骨架；无储能基线 $C^{base}$ 口径；
   储电水价值与阈值律语言};

\draw[ar] (q2.south) -- (q3.north);
\node[lab] at (0.35,-5.43)
  {\textbf{提供}\ 固定场景集（消融与 VoI 共用）；报童临界分位数给出的安全边际；
   紧急购电的结构性认识};

\draw[ar] (q3.south) -- (q4.north);
\node[lab] at (0.35,-7.08)
  {\textbf{提供}\ 多阶段滚动框架与非预期性实现；终端条件处理；
   预报频次对照（VoI）方法};

% ---------------- 输出 ----------------
\node[obox] (o1) at (-5.55,-2.95) {result1.xlsx\\+ 表 1 / 表 2};
\node[obox] (o2) at (-5.55,-4.60) {result2.xlsx\\+ 表 3};
\node[obox] (o3) at (-5.55,-6.25) {result3.xlsx\\+ VoI 结论};
\node[obox] (o4) at (-5.55,-7.90) {result4-2.xlsx\\result4-3.xlsx};

\foreach \q/\o in {q1/o1,q2/o2,q3/o3,q4/o4}{
  \draw[ar, draw=green!45!black] (\q.west) -- (\o.east);
}

\end{tikzpicture}

  \caption{四问的模型依赖关系与复用内容（箭头标注上游向下游提供的模型要素；四问在 2025-02-01 至
            12-31 共 334 天的同一窗口上报告结果）}
  \label{fig:dependency}
\end{figure}

四条上游链路的误差传播与处理见表~\ref{tab:errorprop}。其中最须防范的是第二节与末节指向的同一风险：
问题 1 的输入是\textbf{全年平均日}，问题 2--4 的窗口只有 334 天，两者是\textbf{两套口径}而非同一量的
两个估计，故分别命名、分别报告（表~\ref{tab:twocaliber}），节省率一律定义为“相对同一窗口、同一电价
口径下的无储能基线“。

\begin{table}[htbp]
  \centering
  \caption{四问之间的误差传播路径与本文的处理方式}
  \label{tab:errorprop}
  \footnotesize
  \setlength{\tabcolsep}{4pt}
  \begin{tabularx}{\linewidth}{L{2.15cm} L{3.05cm} Y}
    \toprule
    上游 \(\rightarrow\) 下游 & 传递的内容 & 潜在的误差传播与本文的处理 \\
    \midrule
    附件 \(\rightarrow\) 问题 1
      & 电价、负载、光伏的时标与单位
      & 时标错位一个槽或丢掉 \(0{:}00{+}1\) 一列会让全部合计系统性偏离，而“全天汇总 \(=\sum\) 144 格”
        的自检对整体错位不可见；本文统一按“标签 \(=\) 时段末端”的\textbf{位置读法}对齐，
        断言每天恰 144 槽并保留末列（\S3.2） \\
    问题 1 \(\rightarrow\) 问题 2
      & 模型骨架、\(C^{base}\) 口径、水价值语言
      & 两套口径量级不同（单日 \(48\,052.0466\) 元对 334 天 \(16\,407\,319.63\) 元），混排一次即失去
        可比性；本文各自命名并在表~\ref{tab:twocaliber} 显式声明 \\
    问题 2 \(\rightarrow\) 问题 3
      & 场景生成与削减方法、安全边际
      & 各问各自生成场景会把“场景噪声”误读成“信息结构改进”；本文把场景集\textbf{一次性固定并存盘}，
        问题 3 与全部消融实验读同一份（\S\ref{sec:ablation}），并以 \(K=60\to120\) 作稳健性检验 \\
    问题 3 \(\rightarrow\) 问题 4
      & 多阶段滚动框架、终端条件
      & 电价随机性与光伏误差叠加，分别抽样会丢失相关结构、低估“高价又缺电”的联合尾部；本文用
        “日净需求残差 \(\times\) 日均电价”二维分层取历史代表日生成\textbf{联合}情景，保留
        \(\rho(p,\text{净负荷})\)（实际 \(0.2620\)、情景 \(0.2597\)） \\
    问题 1--4 \(\rightarrow\) 结论
      & 各问的“总费用”与“节省率”
      & 基线不一致时同一“节省率”指向不同分母；四问统一使用 \calref{C-01} 的无储能基线与
        \calref{C-04} 的节省率，并逐处注明电价口径 \\
    \bottomrule
  \end{tabularx}
\end{table}

\begin{table}[htbp]
  \centering
  \caption{无储能基线的两套口径声明（费用单位：元；\calref{C-01} 定义见\chapref{chap:q1}）}
  \label{tab:twocaliber}
  \footnotesize
  \setlength{\tabcolsep}{4pt}
  \begin{tabularx}{\linewidth}{L{2.6cm} C{1.5cm} L{4.6cm} C{1.95cm} C{2.0cm}}
    \toprule
    口径名称 & 适用问题 & 负载与光伏来源 & 电价来源 & 基线数值 \\
    \midrule
    典型日口径 & 问题 1 & 附件 1 的 144 槽典型日（全年逐时段均值） & 附件 1（固定价）
      & \(48\,052.0466\) \\
    窗口口径（固定价） & 问题 2、3 & 附件 2 的当日实际负载与实际光伏，逐日计 & 附件 1（固定价）
      & \(16\,407\,319.63\) \\
    窗口口径（波动价） & 问题 4 & 同上一行 & 附件 4（波动价） & \(17\,026\,782.86\) \\
    \bottomrule
  \end{tabularx}
\end{table}

\subsection{与既有做法的差异}
各问单独选口径、或用元启发式直接求解，会使节省率的分母不可核对且拿不到对偶信息。本文的三点差异：
① 四问共用\textbf{同一物理骨架与同一基线口径}，逐层增加模块，每层贡献可单独消融；
② 全部模型为\textbf{线性规划}并用 HiGHS 精确求解，可输出对偶量 \(\nu_t\) 与阈值型策略的最优性证据；
③ 对题面“是否需要引入其他时刻的预报”给出\textbf{可判定阈值}：同一场景集与同一实际轨迹下，
预报时刻数 \(n=1,2,4\) 的总费用为 \(1\,426.48\)、\(1\,413.16\)、\(1\,393.24\) 万元，
\(n=4\) 已实现 \(85.1\%\) 的可获得信息价值，继续加密最多再省 \(58\,063\) 元（\(0.42\%\)），
故在本题预报精度与惩罚结构下\textbf{不建议继续加密}（\evid{C-Q3-01}）。
